\subsection{局部变化率、链式法则与偏导数}

前述割线斜率的极限说明，导数不是“两个无穷小量的普通商”，而是一个局部变化率。位置对时间的一阶导数是速度，二阶导数是
加速度；电量对时间的导数是电流；弹性体的切线刚度可写成 $\dd F/\dd x$；若物体吸热量为
$Q(T)$，则温度为 $T$ 时的热容 $C(T)$ 为 $\dd Q/\dd T$。这些定义都在问同一件事：系统处于某个状态时，
再作极小改变会得到怎样的即时响应。

复合关系要用链式法则。若 $y=f(u)$、$u=g(x)$，则
\[
 \frac{\dd y}{\dd x}=\frac{\dd y}{\dd u}\frac{\dd u}{\dd x}.
\]
例如势能 $U(x)$ 随位置变化，而位置又随时间变化，就有
$\dd U/\dd t=(\dd U/\dd x)v$。若函数由多个物理量共同决定，则应说明求偏导时保持哪些量不变；
热学中 $(\partial p/\partial V)_T$ 与 $(\partial p/\partial V)_S$ 就代表不同的物理过程。

\subsection{极值判别、物理约束与近似展开}

若函数在定义域内部的 $x_0$ 处可导并取得局部极值，则必有 $f'(x_0)=0$；但驻点不一定是极值，
还要比较左右导数或二阶导数。若 $f'(x_0)=0$ 且 $f''(x_0)>0$，则 $x_0$ 是严格局部极小点；
若 $f'(x_0)=0$ 且 $f''(x_0)<0$，则 $x_0$ 是严格局部极大点。端点和不可导点也可能
给出最大、最小值，故完整求极值必须连同定义域一起检查。物理答案还要满足长度为正、接触不脱离、
速度可达等约束，纯代数驻点未必可实现。

在 $x_0$ 附近，足够光滑的函数可逐阶展开：
\[
 f(x_0+\xi)=f(x_0)+f'(x_0)\xi+
 \frac12f''(x_0)\xi^2+\frac1{3!}f'''(x_0)\xi^3+\cdots .
\]
近似保留到哪一阶取决于所需精度。若 $x_0$ 是势能的平衡位置，一阶项消失，最低非零项通常是二次项，
系统因而近似简谐振动；若二次项也为零，就必须保留更高阶，回复力和振动周期会表现出不同规律。
小角近似 $\sin\theta\simeq\theta$、$\cos\theta\simeq1-\theta^2/2$ 也来自同一展开，使用前要说明
$|\theta|\ll1$ 且角度以弧度计。

对时间求导时还常使用点记号：$\dot X=\dd X/\dd t$，
$\ddot X=\dd^2X/\dd t^2$；对一般自变量仍写撇号或 $\dd/\dd x$。
导数还提供变化率的量级判断。设某量的特征值为 $X$、变化时间为 $T$，则
$\dot X$ 的量级约为 $X/T$；两个机制谁占主导，可比较相应导数项的无量纲比值。竞赛推导中，
这类先估量级、再保留主导项的做法，常比直接处理完整方程更清楚。
